\documentclass[11pt]{article}
\usepackage[margin=0.78in]{geometry}
\usepackage{graphicx}
\usepackage{booktabs}
\usepackage{amsmath}
\usepackage{siunitx}
\usepackage{xcolor}
\usepackage[hidelinks]{hyperref}
\usepackage{microtype}
\definecolor{statusblue}{HTML}{EAF2F8}
\newcommand{\eps}{\ensuremath{\epsilon}}
\newcommand{\CLs}{\ensuremath{CL_s}}
\newcommand{\StageToys}{300}
\newcommand{\StageToyMin}{0}
\newcommand{\StageToyMax}{299}
\newcommand{\ScopeMassRows}{680}
\newcommand{\ToyLimitRows}{204,000}
\newcommand{\TotalWindowRows}{232}
\newcommand{\EmpiricalPResolution}{0.00333}
\newcommand{\CenteredProfileRetries}{2}
\newcommand{\FreeProfileRetries}{1}
\newcommand{\NullProfileRetries}{1}
\newcommand{\StageStatusTitle}{Cumulative 300-toy stage}
\newcommand{\StageStatusText}{This is the planned full cumulative ensemble; all earlier toy IDs are retained exactly.}
\newcommand{\ContinuationText}{The contracted cumulative schedule is complete at 300 toys per mass.}
\newcommand{\GlobalMinMass}{66}
\newcommand{\GlobalGridCount}{232}
\newcommand{\GlobalNeff}{35.3814}
\newcommand{\GlobalLocalP}{\num{0.00284}}
\newcommand{\GlobalSidakP}{\num{0.0959}}
\newcommand{\GlobalSidakZ}{1.305}
\newcommand{\GlobalBonferroniP}{\num{0.66}}
\newcommand{\AllScopesBonferroniP}{\num{0.0861}}
\newcommand{\StageDate}{September 5, 2026}

\title{HPS Gaussian-Process Resonance Search\\
Analysis Note v4.9.12: Expected-Limit Bands and Observed Diagnostics}
\author{Emrys Peets}
\date{\StageDate}

\begin{document}
\maketitle

\begin{center}
\fcolorbox{blue!35!black}{statusblue}{%
\parbox{0.91\linewidth}{\small
\textbf{\StageStatusTitle.} This results-only note uses \StageToys{} background-only
toys per mass (toy IDs \StageToyMin--\StageToyMax). \StageStatusText{} All bands
and empirical limit-tail diagnostics are pointwise and conditional on frozen GP
states.}}
\end{center}

\section{Inputs and construction}

The seven scopes are the exact final-sample curves used in the most recent
Harvard results section: full 2015, full 2016, the optimized-support 2021 10\%
sample, their three pairs, and the all-three shared-$\eps^2$ combination. The
2021 configuration uses the frozen 36--300 MeV GP support and factor-15
resolution-scaled length ceiling. No 2021 1\% or 2016 10\% comparison enters
these results.

At each mass, a latent background vector is drawn from the conditioned
mass-local GP prediction and then Poisson sampled. The same dataset draw is
reused by every scope containing that dataset at a fixed mass and toy ID, so
the standalone and combined comparisons are paired. GP hyperparameters are not
refit inside toys. Each pseudoexperiment is evaluated with the same bounded,
piecewise-asymptotic 90\% \CLs{} construction as the observed v4.9.12 curve.
For numerical stability, fixed-strength nuisance fits minimize a
data-constant-centered Poisson deviance; the fitted likelihood, statistic, and
root acceptance gates are unchanged from the parent solver. During the 100-toy
continuation, an inherited free-strength optimizer status failed at one
deterministic coordinate despite a finite candidate. The contracted solver now
uses the same likelihood-equivalent centering as a retry only when a raw free fit
is unsuccessful. The released ensemble records \CenteredProfileRetries{} centered-profile
retry activations (\FreeProfileRetries{} free-profile and \NullProfileRetries{} null-profile);
every retry still passes the original likelihood-nesting and root gates, and no toy is dropped
or reseeded.

The bands are empirical 2.5\%, 16\%, 50\%, 84\%, and 97.5\% quantiles. Because
different mass points are generated independently, these are pointwise
expected-limit bands rather than a scan-wide ensemble. They do not establish
unconditional coverage, toy-calibrated \CLs{}, a look-elsewhere correction, or
global significance. Full-2016 and 2016-inclusive curves also retain the
previously disclosed cross-process GP-state replay exception.

\section{Results}

\begin{figure}[htbp]
  \centering
  \includegraphics[width=0.99\linewidth]{../figures/final_total_search_window_expected_bands_\StageToys toys.pdf}
  \caption{The ``final'' mass-by-mass maximal-available-dataset observed limit and
  expected bands over the total search window, 19--250 MeV. The scope strip
  states which final datasets enter each interval: 2015 at 19--38 MeV,
  2015+2016 at 39--49 MeV, all three at 50--90 MeV, 2016+2021 at
  91--180 MeV, and 2021 at 181--250 MeV. Scope boundaries are explicit because
  this stitched display is not one invariant likelihood across the whole range.}
  \label{fig:total-window-bands}
\end{figure}

Figure~\ref{fig:total-window-bands} supplies the requested standalone total-window
view. It always selects the combination containing every final dataset available
at that mass; it does not select among scopes according to the observed or expected
limit. Thus the stitching rule is fixed by dataset support rather than by a favorable
fluctuation. The displayed 2021 pieces use the optimized-support 10\% configuration.

\begin{figure}[htbp]
  \centering
  \includegraphics[width=0.98\linewidth]{../figures/all_three_expected_bands_\StageToys toys.pdf}
  \caption{Observed and expected 90\% \CLs{} upper limit for the all-three
  shared-coupling combination. Yellow and green show the central 95\% and 68\%
  pointwise background-only ranges; the dashed curve is the expected median.
  The legend is outside the data region so it does not obscure any curve.}
  \label{fig:all-three-bands}
\end{figure}

Figure~\ref{fig:all-three-bands} is the direct Brazil-band analogue of the
earlier v4.2/v4.5 combination display, now evaluated on the current final
inputs. Local motion in the outer envelope should still be read with the
finite-ensemble granularity of the current \StageToys-toy stage in mind.

\begin{figure}[p]
  \centering
  \includegraphics[width=0.99\linewidth]{../figures/combination_expected_band_panels_\StageToys toys.pdf}
  \caption{Pairwise and all-three shared-$\eps^2$ limits. Each combination has
  its own axis, preventing bands or observed curves from covering another
  result. All four panels use paired constituent pseudo-observations.}
  \label{fig:combination-panels}
\end{figure}

\begin{figure}[p]
  \centering
  \includegraphics[width=0.99\linewidth]{../figures/individual_expected_band_panels_\StageToys toys.pdf}
  \caption{Standalone final-sample limits. The 2021 panel is the optimized
  native-10\% configuration with 36--300 MeV GP support.}
  \label{fig:individual-panels}
\end{figure}

\clearpage
\subsection{Observed-limit and analytic p-values}

Two different diagnostic families are reported and must not be conflated. For an
observed upper limit $U_{\rm obs}(m)$ and toy limits $U_t(m)$ at a fixed mass,
\begin{align}
p_{\rm strong}(m)&=\frac{1}{N}\sum_t
  \mathbf{1}\!\left[U_t(m)\leq U_{\rm obs}(m)\right],\\
p_{\rm weak}(m)&=\frac{1}{N}\sum_t
  \mathbf{1}\!\left[U_t(m)\geq U_{\rm obs}(m)\right],\\
p_{\rm two}(m)&=\min\!\left[1,2\min(p_{\rm strong},p_{\rm weak})\right].
\end{align}
A small $p_{\rm strong}$ marks an unusually tight observed limit, while a small
$p_{\rm weak}$ marks an unusually loose observed limit. The two-sided number is a
symmetric reviewer diagnostic. These are upper-limit ensemble positions, not
discovery p-values and not coverage measurements. The one-sided fractions have
granularity $1/N=\EmpiricalPResolution$, while $p_{\rm two}$ is constructed from
them; zero counts are retained as zero in the CSV and shown at half a count only
for logarithmic plotting.

The analytic $p_0$ is separate: it is the frozen observed curve's one-sided,
fixed-mass, asymptotic profile-likelihood background-only p-value. It is not derived
from these \StageToys{} band toys. No look-elsewhere correction is applied, so neither
the analytic minima nor the empirical tail fractions support a global-significance
claim.

\begin{figure}[p]
  \centering
  \includegraphics[width=0.99\linewidth]{../figures/combination_pvalue_panels_\StageToys toys.pdf}
  \caption{For each shared-$\eps^2$ combination, the three empirical observed-limit
  tail diagnostics and the analytic local discovery $p_0$. A downward triangle at
  half the empirical resolution denotes exactly zero toys in that tail; it is not a
  nonzero p-value estimate. The horizontal dotted line is 0.05 and is included only
  as a visual reference.}
  \label{fig:combination-pvalues}
\end{figure}

\clearpage
\begin{table}[htbp]
\centering
\scriptsize
\setlength{\tabcolsep}{5.1pt}
\begin{tabular}{lrrrrrr}
\toprule
Scope & $m$ [MeV] & $p_{\rm strong}$ & $p_{\rm weak}$ & $p_{\rm two}$ & analytic $p_0$ & $Z_{\rm local}$ \\
\midrule
2015 full & 51 & \num{1} & \num{0} & \num{0} & \num{0.000846} & 3.139 \\
2016 full & 90 & \num{1} & \num{0} & \num{0} & \num{0.000309} & 3.424 \\
2021 10\% & 78 & \num{0.99} & \num{0.01} & \num{0.02} & \num{0.00247} & 2.810 \\
2015 + 2016 & 90 & \num{1} & \num{0} & \num{0} & \num{0.000127} & 3.659 \\
2015 + 2021 & 50 & \num{1} & \num{0} & \num{0} & \num{0.000403} & 3.351 \\
2016 + 2021 & 66 & \num{0.99} & \num{0.01} & \num{0.02} & \num{0.0065} & 2.484 \\
All three & 66 & \num{0.99} & \num{0.01} & \num{0.02} & \num{0.00284} & 2.765 \\
\bottomrule
\end{tabular}
\caption{Fixed-mass diagnostics at each scope's smallest analytic local $p_0$. The one-sided empirical fractions have granularity $1/100=0.010$ and $p_{\rm two}$ is constructed from them; the selected minimum analytic $p_0$ is a compact reporting choice, not a trials-corrected global test.}
\label{tab:pvalue-summary}
\end{table}


Table~\ref{tab:pvalue-summary} gives explicit values at the smallest analytic local
$p_0$ in each scope. Selecting that row makes the compact table scan-aware but does
not calibrate the minimum; interpretation remains fixed-mass/local. Complete values
for all \ScopeMassRows{} scope--mass coordinates are in the machine-readable p-value
ledger.

\subsection{Compact limit summary}

Table~\ref{tab:band-minima} reports, for each scope, the mass at which the
\StageToys-toy expected median is smallest. These minima summarize the plotted
curves; they are not optimization inputs and do not select the analysis card.

\begin{table}[htbp]
\centering
\small
\begin{tabular}{lrrrrr}
\toprule
Scope & $m$ [MeV] & Observed & $-2\sigma$ & Median & $+2\sigma$ \\
\midrule
2015 full & 37 & \num{7.582e-06} & \num{3.520e-06} & \num{5.757e-06} & \num{1.180e-05} \\
2016 full & 73 & \num{1.600e-06} & \num{1.996e-06} & \num{3.467e-06} & \num{5.820e-06} \\
2021 10\% & 102 & \num{2.726e-06} & \num{1.119e-06} & \num{1.903e-06} & \num{4.215e-06} \\
2015 + 2016 & 73 & \num{1.548e-06} & \num{1.948e-06} & \num{3.334e-06} & \num{6.212e-06} \\
2015 + 2021 & 80 & \num{5.005e-06} & \num{1.294e-06} & \num{2.177e-06} & \num{4.815e-06} \\
2016 + 2021 & 102 & \num{1.360e-06} & \num{1.078e-06} & \num{1.752e-06} & \num{3.194e-06} \\
All three & 86 & \num{1.429e-06} & \num{1.023e-06} & \num{1.847e-06} & \num{3.763e-06} \\
\bottomrule
\end{tabular}
\caption{Observed and conditional expected $\epsilon^2$ limits at each scope's minimum expected-median mass. The outer columns are empirical central-95\% endpoints from 50 toys per mass.}
\label{tab:band-minima}
\end{table}


Across all \ScopeMassRows{} scope--mass coordinates, all \ToyLimitRows{} toy
limits were finite and closed at \CLs{}=0.10 within the frozen numerical
tolerance. No latent GP draw used the clipping fallback. The exact per-toy
limits and pseudo-observation hashes are retained in machine-readable ledgers.

\section{Continuation}

The ensemble is cumulative. \ContinuationText{} Atomic per-mass checkpoints make
the continuation resumable, and the contract hash prevents a
changed card, GP-state ledger, observed curve, solver, or protocol from being
silently mixed into the existing toys.

\end{document}
